\printedpage{19}
% head: A. Grothendieck : Dérivateurs, Ch. VII, Édition des Universités de Montpellier II et Paris VII
% foot: Version 22 décembre 2001

\msp{page 29}
But since
\[
\left\{
\begin{array}{l}
\sigma(c)=c(0_\tau),\ \beta(c)=c(1_\tau)\\[1ex]
\text{where } 0_\tau,\ 1_\tau \text{ are the source and target objects of the standard oriented interval } I_\tau\ ,
\end{array}
\right.
\]
and similarly for $\tau'$, $v$ will indeed induce
\[
\underbrace{\sigma(c')}_{=c'(0_{\tau'})} \longrightarrow c(\varphi(0_{\tau'})) \overset{?}{=} \underbrace{\sigma(c)}_{c(0_\tau)},
\]
and likewise for $\beta(c')$ and $\beta(c)$, only if we impose
\[
\boxed{
\left\{
\begin{array}{lcl}
\varphi(0_{\tau'}) & = & 0_\tau\\
\varphi(1_{\tau'}) & = & 1_\tau,
\end{array}
\right.
}
\]
i.e. that the functor $\varphi$ under consideration be \textit{compatible with the source and the target.} This is what visibly had to be imposed in the definition of % [sic] C'est ce visiblement -> C'est visiblement ce, word order
\[
\mathrm{Hom}_{\ul{Ty}}(\tau',\tau)=\underline{\mathrm{Hom}}^{\circled{int}}(I_{\tau'},I_\tau)\quad\text{-}
\]
one must take the functors $\varphi$ respecting the interval structures, i.e. making commutative the diagram
\[
\begin{tikzcd}[column sep=small]
I_{\tau'} \arrow[rr] & & I_\tau \\
& \varepsilon=\{0,1\} \arrow[ul, hook', "i_{\tau'}"] \arrow[ur, hook, "i_\tau"'] &
\end{tikzcd}
\]
defining the orientation of $I_{\tau'}$ and of $I_\tau$ (\footnote[10]{In fact, one must probably restrict the morphisms still further by requiring that they be \textit{increasing} on objects - cf. page 31.}). This rectification made, I have good hope that the category $\ul{Ch}(X)$ we have just constructed is ``the right one''. One could denote by $\ul{CH}(X)$ the larger category, over the category $\ul{Typ}$ having the same objects as $\ul{Ty}$, but more morphisms: namely % [sic] (in footnote 10) qu'il soient -> qu'ils soient, missing plural s
\msp{page 30}
\[
\begin{array}{lcl}
\mathrm{Hom}_{\ul{Typ}}(\tau,\tau') & = & \mathrm{Hom}_{\mathrm{Cat}}(I_\tau,I_{\tau'})\\[0.5ex]
\mathrm{Hom}_{\ul{Ty}}(\tau,\tau') & = & \mathrm{Hom}^{\mathrm{int}}_{\mathrm{Cat}}(I_\tau,I_{\tau'})\\[0.5ex]
 & & \text{(homomorphisms respecting the interval structure).}
\end{array}
\]
We will have
\[
\ul{Ch}(X)\simeq\ul{CH}(X)\times_{\ul{Typ}}\ul{Ty},
\]
\printedpage{20}
% head: A. Grothendieck : Dérivateurs, Ch. VII, Édition des Universités de Montpellier II et Paris VII
% foot: Version 22 décembre 2001
i.e. $\ul{Ch}(X)$ is the \textit{restriction} of the fibered category $\ul{CH}(X)$ over $\ul{Typ}$, to the (non-full) subcategory $\ul{Ty}$ of $\ul{Typ}$.

\textbf{NB} The fibered categories $\ul{CH}(X)$, $\ul{Ch}(X)$ over $\ul{Typ}$, $\ul{Ty}$ depend functorially on $X$, varying in (Cat), as it should be. As is only right, we are interested in them above all when $X$ \textit{is not} in (Cat), i.e. when $X$ is a ``big'' category (such as (Cat) itself). I dispense with making explicit here the pertinent Universe questions.

Two essential tests, to ascertain whether the notion is good
\begin{enumerate}
\item[a)] The theory of the $\mathcal{M}\Sigma^{-1}$.
\item[b)] The composition of paths must work well.
\end{enumerate}

\msp{page 31}
We have the diagram of canonical functors
\[
\begin{tikzcd}
\ul{Ty} \arrow[r, hook] \arrow[d, hook] & \varepsilon\backslash\mathrm{Cat} \arrow[d] \\
\ul{Typ} \arrow[r, hook] & \mathrm{Cat},
\end{tikzcd}
\]
where the horizontal arrows are fully faithful.

\printedpage{21}
% head: A. Grothendieck : Dérivateurs, Ch. VII, Édition des Universités de Montpellier II et Paris VII
% foot: Version 22 décembre 2001

\section*{5\quad Properties of morphisms of intervals}

I have more or less convinced myself that to lend itself to calculations and to simple arguments, one must restrict the morphisms still further in the definition of the category $\ul{Ty}$. I will require that the morphisms $\tau\longrightarrow\tau'$ be functors
\[
w : I_\tau\longrightarrow I_{\tau'}
\]
which are compatible with the \textit{oriented interval} structure of $I_\tau$, $I_{\tau'}$, i.e. that they take source to source, target to target, but moreover that they be \textit{increasing} on objects (for the order relation of progression of paths). This is a restriction which geometrically seems quite natural to me. (It is similar to the analogous restriction, in the pre-categorical definition of morphisms between $\bDelta_i$ \ldots\,) It is this hypothesis that allows one to make a fairly simple structure theory of morphisms of intervals. But for the moment I am going to distinguish ``morphisms'' and ``increasing morphisms''.

\msp{page 32}
Let
\[
f : I_\tau\longrightarrow I_{\tau'}
\]
be a morphism of oriented intervals. Consider the full subcategories, isomorphic to $\Delta_i^+$ or $\Delta_i^-$ as oriented intervals, defined by the critical intervals of $I_\tau$. They form a sequence $\Sigma_\tau$ of subcategories of $I_\tau$
\[
\Sigma_\tau=(\bDelta^{\varepsilon_1}_{l_1}(\tau),\bDelta^{\varepsilon_2}_{l_2}(\tau),\ldots,\bDelta^{\varepsilon_r}_{l_r}(\tau))
\]
of which $I_\tau$ is the amalgamated sum (we will come back to this point further \texttt{[down]}), $r$ being the useful length of $\tau$. We call them the \textit{critical subintervals} of $I_\tau$. Thus $f$ decomposes into a sequence of morphisms
\[
f_i : \bDelta^{\varepsilon_i}_{l_i}\longrightarrow I_{\tau'},\qquad i\in[1,r],
\]
such that
\begin{enumerate}
\item[a)] for $i\in[1,r-1]$, the target of $f_i$ is equal to the source of $f_{i+1}$.

Moreover,
\item[b)] $f_1$ takes source to source, and $f_r$ takes target to target.
\end{enumerate}

Conversely, when we are given $f_i$ satisfying these two conditions, they define a unique morphism of intervals $f$. (And if condition b) is not required, we find a description of the morphisms of categories $I_\tau\longrightarrow I_{\tau'}$.)

\msp{page 33}
We are first going to look at an arbitrary functor
\[
f : \bDelta^{\varepsilon}_i\longrightarrow I_{\tau'}\qquad(\varepsilon\in\{\pm1\})\quad .
\]

\printedpage{22}
% head: A. Grothendieck : Dérivateurs, Ch. VII, Édition des Universités de Montpellier II et Paris VII
% foot: Version 22 décembre 2001
Suppose to fix ideas that $\varepsilon=+1$, i.e. that the first member is a direct critical interval \ldots{} We find that $f$ factors through a critical subinterval of $I_{\tau'}$, say $\bDelta^{\varepsilon'}_j$. It is therefore defined by an increasing map from $\bDelta^{\varepsilon}_i$ to $\bDelta^{\varepsilon'}_j$ (if $\varepsilon$, $\varepsilon'$ \texttt{[are]} of the same sign), a decreasing map in the opposite direction (\footnote[11]{\textbf{Lemma.} \textit{If $I_{\tau'}$ is an interval, every subset $J$ of $I_{\tau'}$ such that for two elements of $J$, there is always an arrow from one to the other, is contained in a critical interval (unique if} $\operatorname{card} J\geq 2$\textit{).}}). But if we require that $f$ be increasing, then it is necessary that \textit{either} $f$ be constant (in which case, if its value is not an endpoint of $I_{\tau'}$, there are exactly \textit{two} critical intervals of $I_{\tau'}$ through which $f$ factors) \textit{or} $f$ nonconstant and $\varepsilon=\varepsilon'$ (and in any case, if $f$ nonconstant, then the critical interval through which $f$ factors is unique). To sum up:

\textbf{Proposition 7.} \textit{Let $I_\tau$, $I_{\tau'}$ be two oriented intervals,}
\[
f : I_\tau\longrightarrow I_{\tau'}
\]
\textit{a functor. Then for every critical interval $J\subseteq I_\tau$ such that $f|J$ is nonconstant, there exists a unique critical interval $J'\subseteq I_{\tau'}$ such that $f|J$ factors through $J'$, i.e. such that $f(J)\subseteq J'$. If $f$ is increasing, then $J$ and $J'$ are of the same sign,}
\msp{page 34}
\textit{i.e. both direct, or both retrograde.}

Thus, if we look at the sequence $s_0,s_1,\ldots,s_n$ of critical vertices of $\tau$, then
\[
t_0=f(s_0),\ t_1=f(s_1),\ \ldots\ ,\ t_n=f(s_n)
\]
form a sequence of vertices of $I_{\tau'}$, such that two consecutive ones are joined by an arrow of $I_{\tau'}$. Thus the full subcategory of $I_{\tau'}$ defined by these $t_i$ is connected, it is therefore a subinterval. If $f$ is increasing, this subinterval is $[t_0,t_n]$. If $f$ is a morphism of intervals, it contains all the critical vertices of $I_{\tau'}$. In all cases, taking $t_i$, $t_j$ $(i<j)$ and applying the above to $f|J$, where $J$ is the subinterval of $I$ with origin $s_i$, endpoint $s_j$, we see that every critical vertex of $I_{\tau'}$ lying between $t_i$ and $t_j$ is among the $t_k$ $(i\leq k\leq j)$. To sum up:

\textbf{Corollary.} \textit{Let $s_i\leq s_j$ be two critical vertices of $I_{\tau'}$. Then every critical vertex of $I_{\tau'}$ lying between $f(s_i)=t_i$ and $f(s_j)=t_j$, is the image of a critical vertex $s_k$ lying between $s_i$ and $s_j$ (i.e. $i\leq k\leq j$). In particular, if $f$ is a morphism of intervals,} % [sic] sommets critiques de I_{\tau'} (first occurrence) -> de I_\tau, the s_i are the critical vertices of \tau
\[
f(I_{\tau\,\red})\supseteq I_{\tau'\,\red}
\]
\textit{(where $I_{\tau\,\red}$ denotes the set of critical vertices of $I_\tau$ or also the reduced subinterval it defines).}

\msp{page 35}
One should be aware that even if $f$ is increasing, one \textit{does not have} in general $f(I_{\tau\,\red})\subseteq I_{\tau'\,\red}$, i.e. $f$ does not necessarily take critical vertices to critical vertices, Minimal example:
\printedpage{23}
% head: A. Grothendieck : Dérivateurs, Ch. VII, Édition des Universités de Montpellier II et Paris VII
% foot: Version 22 décembre 2001

\noindent
\begin{minipage}[c]{0.6\textwidth}
\[
\begin{tikzcd}
\Phi_3^+\ : & s_0 \arrow[r] \arrow[d, mapsto] & s_1 \arrow[d, mapsto] & s_2 \arrow[l] \arrow[r] \arrow[dl, mapsto] & s_3 \arrow[dl, mapsto] \\
\bDelta_2^+\ : & {} \arrow[r] & {} \arrow[r] & {} &
\end{tikzcd}
\]
\end{minipage}\hfill
\begin{minipage}[c]{0.35\textwidth}
$\Phi_3^+$ is reduced, the two critical vertices taken to the noncritical vertex 1.
\end{minipage}

Let us denote by $\mathrm{Cr}(\tau)$ the set of critical intervals of $I_\tau$, in one-to-one correspondence (respecting the total order) with the set of transition arrows of $I_{\tau\,\red}$, % [sic] l'ensembles -> l'ensemble, spurious plural s
\[
\mathrm{Cr}(\tau)\simeq\mathrm{Fl}\,\mathrm{tr}\,(I_{\tau\,\red}).
\]
Let, for a functor
\[
f : I_\tau\longrightarrow I_{\tau'},
\]
$\mathrm{Cr}(\tau,f)$ be the set of critical intervals of $I_\tau$ on which $f$ is not constant. We then have a map
\[
\mathrm{Cr}(\tau,f)\longrightarrow\mathrm{Cr}(\tau') \tag*{$(*)$} % label (*) printed at the left margin
\]
induced by $f$ via proposition 7, associating to each $J\in\mathrm{Cr}(\tau,f)$ the unique $J'\in\mathrm{Cr}(\tau')$ such that $f(J)\subseteq J'$. We have moreover

\textbf{Corollary 2.}
\begin{enumerate}
\item[\textit{a)}] \textit{The image of $(*)$ is a segment of the totally ordered set} $\mathrm{Cr}(\tau')$\textit{, i.e. if two critical intervals $J'$, $J''$ in $I'_{\tau'}$ are in the image, the same is true of all} % [sic] I'_{\tau'} -> I_{\tau'}, spurious prime
\msp{page 36}
\textit{those lying between them.}
\item[\textit{b)}] \textit{This image also contains the critical interval of $I_{\tau'}$ (there is at most one) which contains $f(0_\tau)$ (resp. $f(1_\tau)$) and whose intersection with the segment $[f(0_\tau),f(1_\tau)]$ has cardinality $\geq 2$, as well as all the critical intervals contained in the said segment.}
\item[\textit{c)}] \textit{If $f$ is a morphism of oriented intervals, then the induced map}
\[
\mathrm{Cr}(f):\mathrm{Cr}(\tau,f)\longrightarrow\mathrm{Cr}(\tau')
\]
\textit{is surjective. More precisely, let $J'\subseteq$} $\mathrm{Cr}(\tau')$\textit{, and let $J=f^{-1}(J')$. Then} $\mathrm{Cr}(f)^{-1}(\{J'\})$ \textit{consists of the critical intervals in} $\mathrm{Cr}(\tau,f)$ contained in $J$. \textit{We have moreover,} if $f$ is increasing, % [sic] J' \subseteq Cr(\tau') -> J' \in Cr(\tau'), J' is an element of Cr(\tau')
\[
\mathrm{long}(J')=\sum_{K\in\mathrm{Cr}(f)^{-1}(\{J'\})}\pi_{\mathrm{n}}(K,f),
\]
\textit{where $\pi_{\mathrm{n}}(K,f)$ is the ``part'' (numerical) of $f|K$, i.e. the length of the interval $[f(0_K),f(1_K)]$. (}\textbf{NB} \textit{If $K$ is a critical interval of $I$ such that $K\notin Cr(\tau,f)$, i.e. $f|K$ is constant, we will put $\pi_{\mathrm{n}}(K,f)=0$. Thus the preceding formula also reads}
\[
\mathrm{long}(J')=\sum\limits_{\underbrace{K\in\mathrm{Cr}(I_\tau),\ K\subseteq J}_{\text{(i.e. }K\,\in\,\mathrm{Cr}(f^{-1}(J'))\text{)}}}\pi_{\mathrm{n}}(K,f).\ \textit{)}
\]
\end{enumerate}